\scanpage{128}
\mspage{-III- \handwritten{6} -}
\begin{flushright}n°\handwritten{383}\end{flushright}
COROLLAIRE 7.- \uline{Soient} $C$ \uline{une chambre et} $x\in E$. \uline{Il existe un élément} $y$ \uline{de} $\overline{C}$ \uline{et un seul et un élément} $g\in G$ \uline{tels que} $x=g(y)$.

En effet la réunion des chambres est le complémentaire de $\displaystyle\handwritten{\bigcup}_{H\in\underline{H}}H$, donc est dense dans $E$. Ainsi $x$ est adhérent à une chambre $g(C)$ $(g\in G)$, soit $x=g(y)$ avec $y\in\overline{C}$. L'unicité de $y$ résulte alors du cor.5 au th.2, \S\ 2, n°3.

COROLLAIRE 8.- \uline{Soient} $x\in E$ \uline{et} $G'$ \uline{le sous-groupe de} $G$ \uline{formé des} $g\in G$ \uline{tels que} $g(x)=x$. \uline{Alors} $G'$ \uline{est engendré par les réflexions} $s_H$ $(H\in\underline{H})$ \uline{telles que} $s_H(x)=x$.

On peut supposer $G$ essentiel, et que $x$ est adhérent à la chambre $C$. Alors $x$ appartient à une face $C_J$ $(J\handwritten{\subset\{}1,2,\ldots,m\handwritten{\}})$ de $C$ (avec les notations du \S\ 2, n°3). Pour $g\in G'$, on a $C_J\handwritten{\cap}g(C_J)\handwritten{\ne\varnothing}$, d'où $gG_J=G_J$ (\S\ 2, n°3, th.2), et $G'\handwritten{\subset}G_J$. Comme $G_J$ laisse $C_J$, et donc $x$, invariant, on a $G'=G_J$. Donc $G'$ est engendré par les réflexions $s_i$ pour $i\handwritten{\notin}J$.

2. \uline{Groupes affines engendrés par des réflexions}.

\uline{Dans ce n°}, $G$ \uline{désigne un sous-groupe du groupe affine de} $E$ (Alg.II, 2ème éd., App.II, n°4), \uline{engendré par des réflexions par rapport à des hyperplans affines}. \uline{On suppose que} $G$, \uline{muni de la topologie discrète}, \uline{opère proprement dans} $E$ (Top.Gén.,III, \S\ 4).

A toute application affine $s$, faisons correspondre l'application linéaire ${}^vs$ associée à $s$ (Alg.II, 2ème éd., App.II, n°4). Si $s$ est la réflexion par rapport à un hyperplan $H$, ${}^vs$ est la réflexion (vectorielle) par rapport à la direction ${}^vH$ de $H$ (ibid., n°3). Ainsi $s\handwritten{\rightsquigarrow}{}^vs$ est un homomorphisme de $G$ sur un sous-groupe ${}^vG$ du groupe orthogonal de $E$, engendré par des réflexions.

Soit $W$ le sous-espace essentiel de $E$ pour ${}^vG$, et soit $W'$ l'orthogonal de $W$ dans $E$. Pour tout $x\in E$ et tout $s\in G$, on a $s(x)-x\in W$. Si $W=E$, autrement dit si ${}^vG$ est essentiel, on dit que $G$ est \uline{essentiel}. Dans le cas général, $G$ définit par passage au
% archive sheet number: 127
